\chapter{Requirements}
\label{chapter:requirements}

The purpose of this chapter is to define the requirements necessary to facilitate a GDPR-compliant productivity benchmarking system. They can split into functional requirements, which must be addressed and solved in this thesis. Non-functional requirements define 


The purpose of this chapter is to define the technical and legal requirements necessary to facilitate a GDPR-compliant productivity benchmarking system. Given the sensitive nature of developer activity data sourced from repositories such as GitHub and Jira, the requirements are derived from a dual perspective: the mandatory legal constraints imposed by the General Data Protection Regulation (GDPR) and the Federal Data Protection Act (BDSG), and the functional necessity to maintain high-fidelity metrics for productivity analysis. This chapter establishes the baseline criteria for pseudonymization, data minimization, and auditability that the proposed architecture must satisfy to ensure 'Data Protection by Design' as mandated by Article 25 of the GDPR.

\section{legal requirements}

lorem ipsum

\section{programming requirements}

 abc